% Originally: content/week-10.tex, \section{Bounded C^1 domains ...} (Corsin Week 10)

\section{Bounded \texorpdfstring{$C^1$}{C1} domains and why the definition in the script is so terrible}
\label{sec:bounded_c1_domains}

% Source: Corsin Nick/Class Notes/Week 10.pdf, p. 6
% Def. 14.3
\begin{definition}[Bounded $C^k$ domain]
\label{def:bounded_ck_domain}
Given $k \geq 1$, a bounded open set $\Omega \subset \mathbb{R}^n$ is called a
\newterm{bounded $C^k$ domain} if $M := \partial\Omega$ is an $(n-1)$-dimensional submanifold
of class $C^k$.
\end{definition}

Citing the 2024 lecture notes, Corsin builds up the machinery needed to actually \emph{use}
this definition. To define \newterm{flux} we need the concept of interior and exterior
normal vectors. For $\Omega$ as above, a normal vector $v \in N_p(\partial\Omega)$ is
\newterm{interior} if
\[p + hv \in \Omega \quad \text{for all } h>0 \text{ small},\]
and \newterm{exterior} if
\[p + hv \in \mathbb{R}^n\setminus\Omega \quad \text{for all } h>0 \text{ small}.\]

A \newterm{Gauss map} is a continuous map of unit normal vectors on $\partial\Omega$. There
exists a unique interior and exterior Gauss map; write
\[\nu \in C^0(\partial\Omega, \mathbb{S}^{n-1}) \text{ such that } \nu(p) \text{ is exterior
  normal}\]
for the exterior unit normal map.

\begin{theorem}[Jordan--Brouwer separation theorem]
\label{thm:jordan_brouwer}
If $M \subseteq \mathbb{R}^{n+1}$ is an $n$-dimensional compact and connected submanifold,
then $\mathbb{R}^{n+1}\setminus M$ has exactly two connected components, an
\newterm{inside} and an \newterm{outside}.
\end{theorem}

% Supplement: Sascha Brack/Class Notes/Week_10_Notes_Monday.pdf, pp. 8--11
